% !TEX ROOT = ../ersti.tex
%
% TODO: stefanzen

%\section{Studiengebühren ade!}
\newpage
\mathphyssecnobar{Studiengebühren ade!}%FIXME
Seit dem Sommersemester 2012 sind die Studiengebühren, die im Sommersemester 2007 in ganz BaWü in Höhe von \EUR{500} eingeführt wurden, wieder abgeschafft. Es erfolgte ein vollständiger finanzieller Ausgleich in Höhe von \EUR{280} aus Landesmitteln, genannt "Qualitätssicherungsmittel", oder liebevoll: "QuaSiMi". Es sind nur \EUR{280} statt \EUR{500}, weil durch Ausnahmeregelungen wie die legendäre "Geschwisterregelung" seit 2009 im Durchschnitt nur dieser Betrag pro eingeschriebener Person an Studiengebühren eingenommen wurde.

Die QuaSiMi sind nicht nur zusätzliche Landesmittel für die Hochschule, sondern auch weiterhin zweckgebunden für die Verbesserung der Studienbedingungen. Außerdem wurde im Gesetz \footnote{StuGebAbschG vom 21.12.2011 \url{http://mwk.baden-wuerttemberg.de/fileadmin/pdf/gesetze/StuGebAbschG/GBl-2011_565_Studiengeb\%C3\%BChrenabschaffungsgesetz_21122011.pdf}} festgeschrieben, dass die QuaSiMi nur im Einvernehmen mit einer legitimierten Vertretung der Studierenden ausgegeben werden dürfen. Das heißt, keine Ausgabe kann gegen den Willen der Studis durchgesetzt werden. Das ist immerhin ein Schritt in Richtung Fairness und Gleichberechtigung.

\subsection*{Gebührenfreiheit?}
Weiterhin müsst ihr jedoch als Teil der \EUR{\beitragssumme}, die ihr jedes Semester an die Uni überweisen müsst, einen "Verwaltungskostenbeitrag" von \EUR{\verwaltungsbetrag} bezahlen. Dieser wird den Hochschulen auf ihren Etat angerechnet, sodass das Geld de facto als Ersparnis ans Land Baden-Württemberg geht. Weil das Land nach der nächsten Wahl die Studiengebühren als Mittel gegen die erste Anhebung des Hochschuletats seit 1996 wiederentdecken könnte, möchten wir das Thema Studiengebühren nicht einfach als Teil der hochschulpolitischen Geschichte abhaken, sondern in diesem Kapitel daran erinnern.
%Dieser Absatz sollte dann rein, wenn unten der historische Exkurs nicht hineinkommt!
\iffalse
Wir laden Euch darüberhinaus ein, Euch über die Hintergründe der Studiengebühren zu informieren, beispielsweise auf unserer Fachschaftshomepage \footnote{http://mathphys.fsk.uni-heidelberg.de/studgeb-historie.html}.
\fi
%

\subsection*{Ein paar Argumente gegen Studiengebühren}

Wie ihr sicher bemerkt habt, gehen wir davon aus, dass es gut sei, wenn es keine Studiengebühren gibt. Was veranlasst uns dazu, mal abgesehen von persönlicher Betroffenheit?
\begin{itemize}
\item {Barrieren zwischen Abi und Studium einreißen!}\\Ein Studium ist verglichen mit einer Ausbildung sehr teuer. Mit Studiengebühren wäre es nochmals knapp \EUR{100}/Monat teurer. Das schreckt nachweislich insbesondere AbiturientInnen ab, deren Eltern selbst nicht studiert haben. Außerdem schreckt es prozentual mehr Mädchen als Jungs ab. Kurzum: Menschen, deren Umgebung von ihnen traditionell weniger erwartet dass sie studieren, studieren weniger; gesellschaftliche Ungleichheit, Stratifizierung und überkommene Rollenbilder werden fortgeschrieben.
\item {Aufhebung von finanziell bedingten Nachteilen im Studium!}\\Studis, die knapp über die BAföG-Grenze fallen, sind am härtesten von Studiengebühren betroffen. Denn bei ihnen reicht das Geld der Eltern oft nicht, um Lebensunterhalt und Gebühren zu finanzieren. Sie müssen also jobben gehen und haben somit deutlich weniger Zeit, sich um ihr eigentliches Studium zu kümmern. Gerade in Zeiten der Master"=Zulassungsbeschränkung führt das im Endeffekt zu finanzieller Auslese.
\item {Bildung ist öffentliches Gut!}\\Die Einführung der Studiengebühren wurde von eingefrorenen Haushälten und massiven Kürzungen im Bildungsbereich seitens der Landesregierungen flankiert. Studiengebühren sind keine Notwendigkeit der Lage, sondern Politik. Unserer Auffassung nach ist Bildung jedoch eins der wichtigsten Güter einer Gesellschaft, und das meinen wir nicht im rein ökonomischen Sinn. Gerade die universitäre Bildung bietet (noch, vergleichsweise) viele Freiräume zum eigenständigen Denken. Dessen Bedeutung ist für eine Gesellschaft, die sich demokratisch nennen möchte, nicht zu unterschätzen. Deshalb sollte Bildung in den Landeshaushalten höchste Priorität haben.
\end{itemize}
%\vfill
%\begin{figure}[h]
%\centering{
%    \includegraphics[width=0.7\textwidth]{bilder/studiengebuehren.png}
%}
%\end{figure}
%\vfill

\subsection*{Verteilung der QuaSiMi an der Universität Heidelberg}

Die Studierenden in Heidelberg setzten durch, dass die Studiengebühren ausschließlich zur Verbesserung der Lehre verwendet werden durften. Um ausreichend Präsenz in den Diskussionen zu haben, wurde auch durchgesetzt dass Studierende in den Kommissionen zur Verteilung der Gelder die absolute Mehrheit haben. Die studentischen VertreterInnen der Kommission werden auf Vorschlag der Fachschaft vom Fakultätsrat gewählt. Diese Gremien haben allerdings nur beratende Funktion. Der von der Kommission beschlossene Verwendungsplan wird dem Fakultätsrat zur Entscheidung vorgelegt. Die Budgetverantwortung liegt beim Dekanat.

Diese Struktur wurde bei Einführung der QuaSiMi oftmals übernommen. An einigen Fakultäten, darunter die für Physik \& Astronomie, wurde die Kompetenz der Mittelverwendung an die Studienkommission übergeben und die Studiengebührenkommission aufgelöst.

\subsection*{Verwendung der QuaSiMi}

Der größte Teil des Gelds, etwa 80\%, geht an Deine Fakultät. Das Rektorat verfügt über 20\% der Mittel als zentrale QuaSiMi, wobei vor jeder Ausgabe das Einvernehmen mit den Studierenden in der zentralen QuaSiMi-Kommission hergestellt werden muss. Daraus geht auch Geld an die zentralen Einrichtungen der Universität, also an die Bibliothek, das Rechenzentrum oder den Hochschulsport. Aktuell werden viele Anträge diverser Einrichtungen und Institute an die zentrale QuaSiMi-Kommission entschieden.

Die einzelnen Fächer haben in den letzten Semestern die ihnen zur Verfügung stehenden Gelder nur teilweise ausgegeben, sodass insgesamt noch eine Summe von
%  STIMMT DIE SUMME??
\EUR{1,3\,Mio} auf den einzelnen Konten der Fakultäten liegen. Deshalb versucht das Rektorat immer wieder durchzusetzen, dass mehr Gelder in zentrale Einrichtungen fließen. Allerdings regt sich hier seitens der Fakultäten und der studentischen Vertretung Widerstand.
%Unklar ist auch, wie viel Prozent das Rektorat einbehalten wird.

Die größten Defizite im Bereich Lehre waren an unseren Fakultäten das schlechte Betreuungsverhältnis, vor allem in den Übungsgruppen und Seminaren, sowie ein ungenügendes Serviceangebot. Um diese Defizite auszugleichen wurde und wird mit den Sondermitteln aus Studiengebühren und QuaSiMi insbesondere in folgende Bereiche investiert:

\vspace{5mm}
\textbf{Mathematik und Informatik}
\begin{itemize}
 \item {Mitarbeiterstellen}\\Lehr-/Servicepersonal, Lehraufträge, AssistentInnen
\item {Hilfskraftmittel}\\ TutorInnen und zusätzliche Übungsgruppen
\item {Ausstattung}\\ Computerpools, Hörsäle, Seminarräume
\item {Materialien}\\ Softwarelizenzen, Skripte, Bücher, eBooks\footnote{\url{http://www.ub.uni-heidelberg.de/helios/epubl/eb/Welcome.html}}, Zeitschriften
\end{itemize}

Weiter wurden Mittel zur Modernisierung der mangelhaften Ausstattung der Hörsäle,
Computerpools, sowie der Bibliothek investiert.

\vspace{5mm}
\textbf{Physik}
\begin{itemize}
\item {Mittelbaustellen}\\Medizinerausbildung, Studienberatung, Lehramtsausbildung, Aufbau des IUP (Glossar) Praktikums
\item {AP + FP Versuche (Glossar)}\\Modernisierung bestehender Versuche, neue
Praktika im FP und dem IUP
\item {Öffnungszeiten}\\Erweiterung der Öffnungszeiten in der \gls{KIP}-Bibliothek sowie dem Studierendensekretariat
\item {Hilfskraftmittel}\\ TutorInnen und zusätzliche Übungsgruppen
\item {Materialien}\\Skripte, Bücher, eBooks\footnote{\url{http://www.ub.uni-heidelberg.de/helios/epubl/eb/Welcome.html}}, Zeitschriften
\end{itemize}

Da experimentelle Erfahrung enorm wichtig für einen Physiker ist, wurde und wird
viel Geld in aktuellste Versuche investiert. Eine Bereicherung stellen sicher auch
die Exkursionen nach vielen Kursvorlesungen dar (CERN, GSI\dots).

%VON 2008...
%Eine genauere aktuelle Auflistung findet ihr im aktuellen MathPhys-Info
%(siehe FS-Homepage\footnote{\url{http://mathphys.fsk.uni-heidelberg.de/mpi.html}}).

%Wenn das Erstiinfo zu fett wird, kann man das rausnehmen      ––– wird gemacht ;)
\iffalse
\subsection*{Historischer Exkurs: Geschichte der Studiengebühren -- Geschichte der Sparprogramme}

Im Wintersemester 1970/71 wurden in der Bundesrepublik Deutschland die allgemeinen Studiengebühren, damals "Hörergeld", mit Hilfe von Protesten und Boykotten abgeschafft.

Aber immer wieder gab es Bewegungen in der Politik, die Studiengebühren forderten. Das ging meistens mit einem Bild von Universität einher, das diese nicht als öffentliche Einrichtung zur Verbreitung des gesellschaftlichen Guts Bildung, sondern als Anbieterin von Leistungen auf dem Markt der Ware Bildung auffasst. Entsprechend wenig ist solche Politik bereit, Bildungseinrichtungen eine angemessene finanzielle Ausstattung zukommen zu lassen.

Es wurden Rückmeldegebühren eingeführt um die Haushaltslöcher der Hochschulen zu stopfen. Die Hochschulen mussten in BaWü Anfang der 1990er arge Kürzungen im Haushalt hinnehmen, der 1996 für 10 Jahre eingefroren wurde (bis auf Baumaßnahmen). Diese geplante finanzielle Austrocknung der Hochschulen wurde übrigens 2006 um weitere 10 Jahre verlängert. Gleichzeitig wurde damals laut darüber nachgedacht, Studiengebühren für Langzeitstudierende einzuführen. Diese wurden in verschiedenen Ländern auch durchgesetzt (in BaWü zum WS 1998/99).

2002 wurde von der Kultusministerkonferenz ein allgemeines Studiengebührenverbot festgeschrieben. Langzeitstudiengebühren waren in bestimmten Ausnahmefällen erlaubt.

Gegen dieses Verbot klagten einige Bundesländer vorm Bundesverfassungsgericht, darunter BaWü. Sie führten an, dass der Bund seine Gesetzgebungskompetenz überschritten und in die Länderkompetenz eingegriffen habe.

Am 26.01.2005 fällte das Bundesverfassungsgericht das Urteil, dass ein Verbot allgemeiner Studiengebühren nur gerechtfertigt sei, um „gleichwertige Lebensbedingungen“ in den Ländern zu wahren. In diesem Fall sei aber kein Anlass zu solcher Sorge gegeben. Somit wurde das Gesetz von 2002 gekippt und der Weg für die allgemeinen Studiengebühren geebnet.

Schon im Dezember 2005 wurde in Baden"=Württemberg Minister Frankenbergs Gesetzentwurf zur Einführung allgemeiner Studiengebühren von der Landesregierung beschlossen. Studiengebühren in Höhe von \EUR{500} pro Semester wurden zum Sommersemester 2007 eingeführt. Es gab weiter Proteste gegen die Einführung, Demos und Boykotts wurden organisiert.

Der AK Studiengebühren der \gls{FSK} organisierte einen Boykott und es wurden Klagen beim Verwaltungsgericht eingereicht. Wie an vielen anderen Universitäten in Baden"=Württemberg wurden die Studierenden in Heidelberg von der studentischen Vollversammlung dazu aufgefordert, die \EUR{500} nicht an die Universität, sondern auf ein Treuhandkonto zu überweisen. Man hätte dann unter Bezug auf das zurückgehaltene Geld als Druckmittel mit der Landesregierung Gespräche begonnen. Leider wurde das beschlossene Quorum mit nur ca. 1200 statt 4500 eingegangenen Zahlungen nicht erreicht und der Boykott daher nicht durchgeführt. Die Studiengebühren konnten damit nicht abgeschafft werden.

Auch im Bildungsstreik, der 2009-2011 die Bildungspolitik aufrüttelte, war die Forderung nach Abschaffung der Studiengebühren (als einer Bildungsgebühr von vielen) stets zentral. Darum organisierten viele baden"=württembergische Bildungsstreik-Gruppen, auch in Heidelberg, nach den Erfahrungen aus Hessen und NRW im Hinblick auf die Landtagswahl nochmals Demonstrationen im Januar 2011. Von SPD über Grüne bis Linkspartei fand sich im Wahlprogramm zur Landtagswahl 2011 denn auch die Absichtserklärung wieder, die Studiengebühren abzuschaffen.

Zum Sommersemester 2012 hat nun tatsächlich die grün-rote Landesregierung ihr Wahlversprechen eingelöst: Die Studiengebühren wurden, bei voller Kompensation aus Landesmitteln, abgeschafft!

%Von Gebühren befreit bist du auch schon in diesem Semester, wenn du ein Kind unter 14 Jahren hast, ein Praxissemester absolvierst, ein Stipendium erhältst, du mindestens zwei (auch Stief"=/Halb"=) Geschwister hast die nicht diese baden"=württembergische Geschwisterregelung in Anspruch nehmen, oder du unter einer studienerschwerenden Behinderung leidest. Eine Befreiung aufgrund außerordentlicher Studienleistungen ist gesetzlich möglich, wird aber von den Fakultäten für Mathematik und Physik abgelehnt und darum nicht gewährt.
\fi
